\section{\texorpdfstring{$\mathcal M(x,y;B)$ 的\termSplittingEnd{}}{M(x,y;B) 的分裂端}}
\label{sec:splitting-ends}

本节将精确描述紧拉格朗日子流形对 $(L_0,L_1)$ 的\personFloer{}流线模空间
$\mathcal M(x,y;B)$ 的\termSplittingEnd{}。另一类端，即\termBubblingEnd{}，将在下一节
描述。

对 $i=0,1$，以
\[
  e_i:\mathfrak j_\omega\longrightarrow\mathcal J_\omega
\]
表示取值映射 $J\mapsto J(i)$。定义
\stfhIndex{I011}{00-j-j0-j1}{$\mathfrak j_{(J_0,J_1)}$}
\[
  \mathfrak j_{(J_0,J_1)}
  =
  \left\{
    J\in\mathcal P([0,1],\mathcal J_\omega)
    \ \middle|\
    e_i(J)=J_i,\quad i=0,1
  \right\},
\]
也就是 $\mathcal J_\omega$ 中分别以 $J_0,J_1$ 为固定端点的道路集合。与之相伴的\personFloer{}模空间为\footnote{译者注：原书所列集合没有写出任何含 $p,q$ 的条件，因而右端实际不依赖
$p,q$。译文依第 13.9 节已经给出的 $\widetilde{\mathcal M}(p,q;J)$ 定义，补出
$u(-\infty)=p$ 与 $u(+\infty)=q$，使当前记号有定义。}
\[
  \begin{aligned}
  \widetilde{\mathcal M}(p,q;J;B)
  =\bigl\{
    u:\Theta\to M\ \bigm|\ &\overline\partial_Ju=0,\quad
    u(\tau,0)\in L_0,\quad u(\tau,1)\in L_1,\\
  &u(-\infty)=p,\quad u(+\infty)=q,\quad
    E_J(u)<\infty,\quad [u]=B
  \bigr\}.
  \end{aligned}
\]
当 $p\neq q$，或者 $p=q$ 且 $0\neq B\in\pi_2(p,p)$ 时，还定义其商
\[
  \mathcal M(p,q;J;B)
  =
  \widetilde{\mathcal M}(p,q;J;B)/\mathbb R.
\]
（其余情形并不稳定，故予以省略。）除非为避免混淆而必须指明 $J$，下文将在
$\widetilde{\mathcal M}(p,q;J;B)$ 与 $\mathcal M(p,q;J;B)$ 的记号中省略 $J$。

一般说来，商空间 $\mathcal M(p,q;B)$ 并不紧。先注意，$\mathbb R$-作用保持同伦类，
因而也保持辛面积。我们将在模去 $\tau$-平移的意义下，分析给定序列
$u_i\in\widetilde{\mathcal M}(p,q;B)$ 紧性失效的方式。

从下面的引理开始。

\begin{lemma}
\label{lem:positive-energy-gap}
设 $L_0,L_1$ 横截相交。则存在只依赖于对子 $(L_0,L_1)$ 与 $J$ 的常数
\[
  \sigma=\sigma_{(L_0,L_1;J)}>0,
\]
使得对方程
\[
  \frac{\partial u}{\partial\tau}
  +J_t\frac{\partial u}{\partial t}=0,
  \qquad
  u(\tau,0)\in L_0,
  \quad
  u(\tau,1)\in L_1
\]
的任意\termNonStationary{}解 $u$，都有
\begin{equation}
  E_J(u)\geq\sigma.
\label{eq:positive-energy-gap}
\end{equation}
\end{lemma}

\begin{proof}
用反证法。假设存在一列\termNonStationary{}解 $u_i$，使得
$E_J(u_i)\to0$（$i\to\infty$）。这尤其排除了冒泡，故得到一致导数界
\[
  |du_i|_\infty<C.
\]
采用引理 14.1.4 证明中的同一论证并抽取子列，仍将其记为 $u_i$，使得
\begin{equation}
  \lim_{i\to\infty}|du_i|_\infty=0.
\label{eq:vanishing-derivative-sequence}
\end{equation}

我们断言，存在邻域 $U\supset L_0\cap L_1$，使得对所有充分大的 $i$ 都有
$\operatorname{Image}u_i\subset U$。

否则，可以找到子列和点 $(\tau_i,t_i)$，使得
\[
  u_i(\tau_i,t_i)\in M\setminus U.
\]
沿 $\tau$-方向平移 $u_i$ 后可设 $\tau_i=0$，并可设
\[
  t_i\longrightarrow t_\infty\in[0,1],
  \qquad
  u_i(0,t_i)\longrightarrow x_\infty\in M\setminus U.
\]
利用~\eqref{eq:vanishing-derivative-sequence}可证 $u_i$ 收敛到一条\personFloer{}流线
$u_\infty$，且 $E_J(u_\infty)=0$。于是 $u_\infty$ 必为常值映射，其像必包含于
$L_0\cap L_1$。另一方面，该收敛又蕴含
$u_\infty(0,t_\infty)=x_\infty\in M\setminus U$。二者显然矛盾，故引理得证。
\end{proof}

接着回忆，引理 14.1.3 蕴含：当 $\tau\to\pm\infty$ 时，$u(\tau)$ 分别收敛到
某个 $p_\pm\in L_0\cap L_1$。\footnote{译者注：严格说，引理 14.1.3 只给出两端各有
一列趋于无穷的平移参数，使相应道路收敛到交点，并未单独证明整个端部
收敛。此处还需结合引理 14.1.4 的导数衰减、横截性给出的交点离散性及标准
紧性论证；原书将这些省略而全部归于引理 14.1.3。}特别地，可将 $u$ 写成
\[
  u(\tau,t)=\exp_{p_\pm}\bigl(\xi_\pm(\tau,t)\bigr),
\]
其中
\begin{align}
  \xi_+:(R_+,\infty)\times[0,1]&\longrightarrow T_{p_+}M,
  \label{eq:positive-end-exponential-coordinate}\\
  \xi_-:(-\infty,R_-)\times[0,1]&\longrightarrow T_{p_-}M.
  \label{eq:negative-end-exponential-coordinate}
\end{align}
这里的指数映射 $\exp_{p_\pm}$ 是相对于任意使 $L_0,L_1$ 在 $p_\pm$ 附近成为\emph{\termTotallyGeodesic{}}子流形的\personRiemann{}度量而言的。$R_\pm$ \emph{依赖于 $u$，并可随
$u$ 变化}。

引理 14.1.3 与命题 14.1.5 于是蕴含：

\begin{proposition}
\label{prop:asymptotic-coordinate-decay}
设 $L_0,L_1$ 横截，并令 $\xi_\pm$ 与 $R_\pm>0$ 如上。\footnote{译者注：原书一方面声明 $R_\pm>0$，另一方面把负端坐标域和估计范围写成
$(-\infty,R_-)$，而通常以正数 $R_-$ 记负端截断时应写 $(-\infty,-R_-)$。
也可能原意是允许 $R_-<0$。两种补法无法由本段唯一裁定，译文保留原式。}则对每个
$k\in\mathbb N$，存在 $C_k>0$ 和只分别依赖于 $p_\pm$ 的
$\delta_\pm=\delta_{p_\pm}>0$，使得对每个给定的 $k\geq0$，
\begin{align}
  |\nabla^k\xi_+|(\tau,t)
  &\leq C_ke^{-\delta_+|\tau|},
  &&R_+<\tau<\infty,
  \label{eq:positive-end-coordinate-decay}\\
  |\nabla^k\xi_-|(\tau,t)
  &\leq C_ke^{-\delta_-|\tau|},
  &&-\infty<\tau<R_-.
  \label{eq:negative-end-coordinate-decay}
\end{align}
\end{proposition}

现在定义 $\widetilde{\mathcal M}(p,q;B)$ 上的一个拓扑；
$\mathcal M(p,q;B)$ 上的拓扑则定义为由此前者诱导的\termQuotientTopology{}。
\stfhIndex{I119}{Jubunengliang}{局部能量（local energy）}
为在实际使用中更有效地研究这种收敛，引入如下\emph{\termLocalEnergy{}}会很有帮助。

\begin{definition}[\termLocalEnergy{}]
\label{def:local-energy}
定义 $u$ 在 $[R_1,R_2]\times[0,1]$ 上的\termLocalEnergy{}为
\[
  E_{J,[R_1,R_2]}(u)
  =
  \int_{R_1}^{R_2}\int_0^1|du|_J^2\,dt\,d\tau.
\]
\end{definition}

由命题 14.2.2 和\termBootstrapArgument{}，得到下面这个更实用的收敛命题。

\begin{proposition}
\label{prop:fine-convergence-criterion}
设序列 $u_i\in\widetilde{\mathcal M}(p,q;B)$ 满足以下条件：
\begin{enumerate}
  \item $u_i$ 在紧集上一致地依 $C^1$ 拓扑收敛到某个映射
  $u\in\widetilde{\mathcal M}(L_0,L_1;J)$；
  \item 局部能量收敛到 $E(u_i)=\omega(B)$，即\footnote{译者注：原书的陈列式只有
  $\lim_{R\to\infty}E_{J,[-R,R]}(u_i)=\omega(B)$，留下自由指标 $i$；而这对每个
  固定的有限能量 $u_i$ 都自动成立，不能排除能量随 $i$ 逃向无穷远。译文依本命题
  的结论及命题 14.2.7 所用的极限次序补入 $i\to\infty$ 的极限。}
  \[
    \lim_{R\to\infty}\lim_{i\to\infty}
    E_{J,[-R,R]}(u_i)=\omega(B).
  \]
\end{enumerate}
则局部极限 $u$ 属于 $\widetilde{\mathcal M}(p,q;B)$，且 $u_i$ 在\termFineCinfinityTopology{}中收敛到 $u$。
\end{proposition}

\begin{exercise}
\label{ex:fine-convergence-criterion}
证明这个命题。
\end{exercise}

依\termQuotientTopology{}的定义，序列 $u_j$（更准确地说是它们的等价类 $[u_j]$）
在 $\mathcal M(p,q;B)$ 中收敛，当且仅当存在序列 $\tau_j\in\mathbb R$，使平移序列
\[
  u_j\circ\tau_j
  :=u_j(\mathord\cdot-\tau_j,\mathord\cdot)
\]
满足上面的两个性质。
\footnote{译者注：原书在此只写未定义的复合 $u_j\circ\tau_j$；译文依本节稍后的
命题 14.2.6，把它明确定义为定义域的 $\tau$-平移。}

这个命题揭示了给定序列何以会收敛失效。受它启发，我们依照\personFloer{} (Fl88b)
引入\emph{\termCuspTrajectory{}}与\emph{\termBrokenCuspTrajectory{}}；它们将用来紧化
$\mathcal M(p,q;B)$。

由于我们已经说明局部拓扑中的收敛失效来自已研究过的冒泡，以下假设已有一致界
\begin{equation}
  |du_i(\tau,t)|<C
\label{eq:uniform-derivative-bound-splitting}
\end{equation}
对所有 $(\tau,t)\in\mathbb R\times[0,1]$ 与 $i\in\mathbb N$ 成立。一般以
$I\subset\mathbb N$ 表示 $\mathbb N$ 的一个子列。下面的命题是\termAscoliArzelaTheorem{}的直接推论。

\begin{proposition}
\label{prop:translated-local-subsequence}
设序列 $u_i\in\widetilde{\mathcal M}(p,q;B)$ 对某个 $C>0$ 满足
\eqref{eq:uniform-derivative-bound-splitting}。则对任意给定序列
$\{\tau_i\}\subset\mathbb R$，都存在 $\mathbb N$ 的一个子列
$I=\{i_j\}_{j\in\mathbb N}$ 和映射
$u\in\widetilde{\mathcal M}(L_0,L_1;J)$，使得平移序列
\[
  u_{i_j}\circ\tau_{i_j}
  =u_{i_j}(\mathord\cdot-\tau_{i_j},\mathord\cdot)
\]
在 $\mathbb R\times[0,1]$ 上局部一致收敛到 $u$。
\end{proposition}

把此命题中与 $I$ 相伴的序列记为\footnote{译者注：原书把右端写成
$\{\tau_{i_j}\}_{j\in I}$，但 $j$ 是枚举子列
$I=\{i_j\}_{j\in\mathbb N}$ 的指标，并不是 $I$ 的元素。译文补正为
$j\in\mathbb N$。}
\[
  \tau_I=\{\tau_{i_j}\}_{j\in\mathbb N}
\]
并把相应的局部极限记为 $u=u_{(I,\tau_I)}$。下面给出局部能量的收敛结论。

\begin{proposition}
\label{prop:local-energy-of-recovered-limit}
设 $u_{i_j}$ 是如上选取的子列，$u_{(I,\tau_I)}$ 是其局部极限。则\footnote{译者注：原书在公式 (14.2.25) 及证明第一行中都对未平移的 $u_{i_j}$
取局部能量，但命题 14.2.6 的局部收敛对象是
$u_{i_j}\circ\tau_{i_j}$。译文补回平移；否则证明中的第一步并不由局部收敛推出。}
\begin{equation}
  \limsup_{R\to\infty}\lim_{j\to\infty}
  \left|
    E_{J,[-R,R]}(u_{i_j}\circ\tau_{i_j})
    -E_J\bigl(u_{(I,\tau_I)}\bigr)
  \right|
  =0.
\label{eq:recovered-local-energy-convergence}
\end{equation}
\end{proposition}

\begin{proof}
由局部收敛，对每个 $R>0$，当 $j\to\infty$ 时有
\[
  E_{J,[-R,R]}(u_{i_j}\circ\tau_{i_j})
  \longrightarrow
  E_{J,[-R,R]}\bigl(u_{(I,\tau_I)}\bigr).
\]
由 $E_J(u_{(I,\tau_I)})<\infty$\footnote{译者注：原书把这里的有限性条件排成
$E_J(u_{(I,\tau_I)})<0$，与能量非负性矛盾。译文补正为 $<\infty$。}，\termLebesgueDominatedConvergenceTheorem{}显然蕴含
\[
  \lim_{R\to\infty}
  E_{J,[-R,R]}\bigl(u_{(I,\tau_I)}\bigr)
  =E_J\bigl(u_{(I,\tau_I)}\bigr).
\]
合并二式即得~\eqref{eq:recovered-local-energy-convergence}。
\end{proof}

下面这个定义会很有用。

\stfhIndex{I121}{Jubujixian}{局部极限（local limit）}
\stfhIndex{I166}{Huifudui}{恢复对（recovering pair）}
\begin{definition}
\label{def:recovering-pair}
设 $\{u_i\}_{i\in\mathbb N}$ 是 $\widetilde{\mathcal M}(p,q;B)$ 中的序列。
把命题 14.2.6 所给的对子 $(I,\tau_I)$ 称为 $\{u_i\}$ 的一个\termRecoveringPair{}。若相应的局部极限 $u_{(I,\tau_I)}$ 不是定常的，则称这个\termRecoveringPair{}是\termNonStationary{}的。也把任何这样的极限
\[
  u_{(I,\tau_I)}\in\widetilde{\mathcal M}(L_0,L_1;J)
\]
称为序列 $u_i$ 的一个\termLocalLimit{}（与\termRecoveringPair{}
$(I,\tau_I)$ 相伴）。
\end{definition}

对给定序列 $\{u_i\}$，收集其所有\termNonStationary{}的\termRecoveringPair{}与\termLocalLimit{}，并分别记为\footnote{译者注：原书第一个集合的竖线后缺少主语，只写
“is a non-stationary recovering pair”；第二个集合又把已经带参数的
$\operatorname{Rcov}(\{u_i\})$ 简写成 $\operatorname{Rcov}$。译文补齐主语并统一
参数。}
\[
  \begin{aligned}
  \operatorname{Rcov}(\{u_i\})
  &=\left\{
      (I,\tau_I)
      \ \middle|\
      (I,\tau_I)\text{ 是 }\{u_i\}\text{ 的非定常恢复对}
    \right\},\\
  \widetilde{\operatorname{Loc}}(\{u_i\})
  &=\left\{
      u\in\widetilde{\mathcal M}(L_0,L_1)
      \ \middle|\
      \begin{gathered}
      u=u_{(I,\tau_I)}\text{ 对某个}\vphantom{\big|}\\
      (I,\tau_I)\in\operatorname{Rcov}(\{u_i\})
      \end{gathered}
    \right\}.
  \end{aligned}
\]
\begin{definition}
\label{def:translation-equivalent-local-limits}
设 $v,v'\in\widetilde{\operatorname{Loc}}(\{u_i\})$ 是 $\{u_i\}$ 的两个\termNonStationary{}的\termLocalLimit{}。若存在 $\tau_0\in\mathbb R$ 使
\[
  v'=v\circ\tau_0,
\]
则称 $v\sim v'$。
\end{definition}

以
\[
  \operatorname{Loc}(\{u_i\})
  =
  \widetilde{\operatorname{Loc}}(\{u_i\})/\mathord\sim
\]
表示等价类集合，并以 $[v]$ 表示与
$v\in\widetilde{\operatorname{Loc}}(\{u_i\})$ 相伴的等价类。
\footnote{译者注：原书在这里写 $v\in\operatorname{Loc}(\{u_i\})$，但 $[v]$
应由商映射作用于商前集合的元素得到。译文补正为
$v\in\widetilde{\operatorname{Loc}}(\{u_i\})$。}
只要含义清楚，下文也把 $[v]$ 本身称为\termLocalLimit{}，并将它记为 $v$。

还需要下面的定义。

\stfhIndex{I052}{Jubujixianlian}{局部极限链（chain of local limits）}
\begin{definition}
\label{def:chain-of-local-limits}
设 $\{u_i\}\subset\widetilde{\mathcal M}(p,q;B)$。所谓 $\{u_i\}$ 的一条\termLocalLimitChain{}，是一列\termLocalLimit{}
\[
  ([v_1],[v_2],\ldots,[v_n])
\]
满足：存在一个子列 $I$ 以及序列
\[
  \tau_j^1<\cdots<\tau_j^n,
  \qquad j\in I,
\]
使得
\[
  v_\ell=u_{(I,\tau_I^\ell)}.
\]
记
\[
  \vec\tau_I
  =
  \left\{(\tau_j^1,\ldots,\tau_j^n)\right\}_{j\in I}.
\]
\end{definition}

以 $\widetilde{\operatorname{Ch}}(\{u_i\})$ 表示 $\{u_i\}$ 的所有\termLocalLimitChain{}组成的集合。若链 $(v_1,\ldots,v_n)$ 与定义 14.2.10
所述子列 $(I,\vec\tau_I)$ 相伴，则记
\[
  C_{(I,\vec\tau_I)}=(v_1,\ldots,v_n).
\]

\stfhIndex{I183}{Zilian}{子链（subchain）}
\begin{definition}
\label{def:subchain-equivalence}
设 $C,C'$ 是 $\{u_i\}$ 的\termLocalLimitChain{}。若
\[
  C=C_{(I,\vec\tau_I)}=(v_1,\ldots,v_n),
\]
且 $I$ 有一个子列 $I'$，使得
\[
  C'=C_{(I',\vec\tau_{I'})}=(v'_1,\ldots,v'_{n'})
\]
并且
\[
  \{v_1,\ldots,v_n\}\subset\{v'_1,\ldots,v'_{n'}\},
\]
则称 $C$ 是 $C'$ 的一条\termSubchain{}。若 $C$ 是 $C'$ 的\termSubchain{}，且
$C'$ 也是 $C$ 的\termSubchain{}，则称两条链等价。
\end{definition}

以 $\operatorname{Ch}(\{u_i\})$ 表示链的等价类集合，并以 $[C]$ 表示 $C$ 的
等价类。现在定义偏序 $[C]\leq[C']$：当且仅当 $C$ 是 $C'$ 的\termSubchain{}。

\begin{lemma}
\label{lem:local-limit-chain-energy-bound}
设 $C=(v_1,\ldots,v_n)$ 是一条长度有限的\termLocalLimitChain{}。则
\begin{equation}
  \sum_{\ell=1}^n\omega([v_\ell])\leq\omega(B).
\label{eq:local-limit-chain-area-bound}
\end{equation}
特别地，链长 $n$ 满足
\begin{equation}
  n\leq
  \left[\frac{\omega(B)}{\sigma}\right]+1.
\label{eq:local-limit-chain-length-bound}
\end{equation}
\end{lemma}

利用这个引理，下面证明\termMaximalChain{}的存在性。

\begin{proposition}
\label{prop:existence-maximal-local-limit-chain}
假设 $p\neq q$，或者 $p=q$ 且 $B\neq0$。设
$u_i\in\widetilde{\mathcal M}(p,q;B)$ 是满足一致估计
\[
  \|du_i\|_\infty<C
\]
的序列。则 $\operatorname{Ch}(\{u_i\})$ 中存在一条\termMaximalChain{}，其长度不
超过 $[\omega(B)/\sigma]+1$。
\footnote{译者注：原书在本命题、命题 14.2.16 和定理 14.2.17 中都把 $u_i$
写成商空间 $\mathcal M(p,q;B)$ 的元素，却同时对 $u_i$ 求导、平移并取逐点值；这些
操作作用于参数化映射而非其平移等价类。译文三处统一补回波浪号，写成
$u_i\in\widetilde{\mathcal M}(p,q;B)$。}
\end{proposition}

\begin{proof}
在给定假设下，$u_i$ 的一致 $C^1$ 界蕴含
$\operatorname{Rcov}(\{u_i\})$ 非空。取一个\termRecoveringPair{}
$(I,\tau_I)$，并令 $u_{(I,\tau_I)}$ 为其\termLocalLimit{}。若
\[
  \omega\bigl(u_{(I,\tau_I)}\bigr)=\omega(B),
\]
则命题 14.2.4 蕴含单元素链
$([u_{(I,\tau_I)}])$ 已经是\termMaximalChain{}。
\footnote{译者注：原书把这条单元素链写成 $\{(I,\tau_I)\}$，其元素是恢复对，
而定义 14.2.10 要求链的元素是局部极限类。译文改写为
$([u_{(I,\tau_I)}])$。}
因此，以下假设
\begin{equation}
  \omega\bigl(u_{(I,\tau_I)}\bigr)<\omega(B).
\label{eq:recovered-energy-deficit}
\end{equation}

由~\eqref{eq:recovered-local-energy-convergence}，对任意 $\epsilon>0$，存在
$R_\epsilon>0$\footnote{译者注：原书在公式 (14.2.29) 的能量区间中固定写
$R_\epsilon$，式后却声称该式对所有 $R\geq R_\epsilon$ 成立，而且没有量化 $j$。
译文依公式 (14.2.25) 把区间变量统一为 $R$，并补明 $j$ 充分大。}，使得当
$R\geq R_\epsilon$ 且 $j$ 充分大时，
\begin{equation}
  \left|
    \omega\bigl(u_{(I,\tau_I)}\bigr)
    -E_{J,[-R,R]}
      \bigl(u_{i_j}(\mathord\cdot-\tau_{i_j},\mathord\cdot)\bigr)
  \right|
  \leq\epsilon.
\label{eq:recovered-energy-window}
\end{equation}
固定 $\epsilon>0$，使得
\[
  \epsilon<
  \frac{\omega(B)-\omega(u_{(I,\tau_I)})}{3}.
\]

记
\[
  u_{i_j,\tau_{i_j}}(\tau,t)
  :=u_{i_j}(\tau-\tau_{i_j},t).
\]
定义 $\mu_{\pm,j}>R$ 为分别落在
$-\infty<-\mu_{-,j}<-R$ 与 $R<\mu_{+,j}<\infty$ 中的唯一数，使得
\begin{equation}
  \begin{aligned}
  E_{J,(-\infty,-\mu_{-,j}]}
    (u_{i_j,\tau_{i_j}})
  &=\frac12E_{J,(-\infty,-R]}
    (u_{i_j,\tau_{i_j}}),\\
  E_{J,[\mu_{+,j},\infty)}
    (u_{i_j,\tau_{i_j}})
  &=\frac12E_{J,[R,\infty)}
    (u_{i_j,\tau_{i_j}}).
  \end{aligned}
\label{eq:half-tail-energy-cutoffs}
\end{equation}

\begin{lemma}
\label{lem:tail-cutoff-divergence}
有
\[
  \max\{\mu_{+,j}-R,\mu_{-,j}-R\}\longrightarrow\infty.
\]
\end{lemma}

\begin{proof}
否则可再取一个仍记为 $j$ 的子列，使得对某个 $C>0$ 和所有 $j$ 都有
\[
  \mu_{\pm,j}\leq R+C.
\]
一方面，\eqref{eq:recovered-energy-window} 蕴含
\[
  \begin{aligned}
  &E_{J,(-\infty,-R-C]}(u_{i_j,\tau_{i_j}})
   +E_{J,[R+C,\infty)}(u_{i_j,\tau_{i_j}})\\
  &\qquad
   =E_J(u_{i_j,\tau_{i_j}})
    -E_{J,[-R-C,R+C]}(u_{i_j,\tau_{i_j}})\\
  &\qquad
   \geq\omega(B)-\omega(u_{(I,\tau_I)})-\epsilon.
  \end{aligned}
\]
另一方面\footnote{译者注：原书在这两串估计中突然改用未平移的 $u_{i_j}$，又把由
公式 (14.2.30) 定义的截点 $-\mu_{-,j}$、$\mu_{+,j}$ 写成
$-\mu_{-,j}-R$、$\mu_{+,j}+R$。这些写法与 (14.2.30) 不能相接。译文统一回
$u_{i_j,\tau_{i_j}}$，并依集合包含关系和 (14.2.30) 使用原截点。}，
由~\eqref{eq:half-tail-energy-cutoffs}，
\[
  \begin{aligned}
  &E_{J,(-\infty,-R-C]}(u_{i_j,\tau_{i_j}})
   +E_{J,[R+C,\infty)}(u_{i_j,\tau_{i_j}})\\
  &\quad\leq
   E_{J,(-\infty,-\mu_{-,j}]}(u_{i_j,\tau_{i_j}})
   +E_{J,[\mu_{+,j},\infty)}(u_{i_j,\tau_{i_j}})\\
  &\quad=
   \frac12\left(
     E_{J,(-\infty,-R]}(u_{i_j,\tau_{i_j}})
     +E_{J,[R,\infty)}(u_{i_j,\tau_{i_j}})
   \right)\\
  &\quad\leq
   \frac12\left(
     \omega(B)-\omega(u_{(I,\tau_I)})+\epsilon
   \right).
  \end{aligned}
\]
合并二式，得到
\[
  \omega(B)-\omega(u_{(I,\tau_I)})-\epsilon
  \leq
  \frac12\left(
    \omega(B)-\omega(u_{(I,\tau_I)})+\epsilon
  \right),
\]
从而
\[
  \omega(B)-\omega(u_{(I,\tau_I)})\leq3\epsilon.
\]
这与 $\epsilon<[\omega(B)-\omega(u_{(I,\tau_I)})]/3$ 矛盾，故引理得证。
\end{proof}

进一步，由~\eqref{eq:recovered-energy-window} 得\footnote{译者注：原书先由
公式 (14.2.30) 把两个远尾的能量之和写成两端总尾能量的一半，继而却断言这个
“一半”仍不小于
$\omega(B)-\omega(u_{(I,\tau_I)})-\epsilon$；公式 (14.2.29) 只能给出两端总尾能量
本身的这一估计。下一句又据此断言某一端不小于能量差的 $2/3$，同样不能推出。
译文先用 (14.2.29) 估计两端总尾能量，再选取至少承载其一半的一端，最后用
(14.2.30) 二分该端；后文所需的正下界相应为能量差的 $1/6$。}
\[
  \begin{aligned}
  &E_{J,(-\infty,-R]}(u_{i_j,\tau_{i_j}})
    +E_{J,[R,\infty)}(u_{i_j,\tau_{i_j}})\\
  &\quad=
    \omega(B)-E_{J,[-R,R]}(u_{i_j,\tau_{i_j}})\\
  &\quad\geq
    \omega(B)-\omega(u_{(I,\tau_I)})-\epsilon\\
  &\quad>
    \frac23\left(
      \omega(B)-\omega(u_{(I,\tau_I)})
    \right).
  \end{aligned}
\]
因而，在必要时再取 $j$ 的子列后，两个量
\[
  E_{J,(-\infty,-R]}(u_{i_j,\tau_{i_j}}),
  \qquad
  E_{J,[R,\infty)}(u_{i_j,\tau_{i_j}})
\]
中至少有一个不小于
\[
  \frac12\left(
    \omega(B)-\omega(u_{(I,\tau_I)})-\epsilon
  \right)
  >
  \frac13\left(
    \omega(B)-\omega(u_{(I,\tau_I)})
  \right)>0.
\]
由 $u_{i_j,\tau_{i_j}}$ 的局部收敛容易看出，这必发生在满足
$\mu_{*,j}-R\to\infty$ 的某个 $*=\pm$ 上。令 $\mu_j$ 为相应的
$\mu_{*,j}$ 之一，使 $\lim_{j\to\infty}\mu_{*,j}=\infty$。不失一般性，可设
$*=-$，因而 $\mu_j=\mu_{-,j}$。

现在取一个子列 $\{i_k\}$，考虑平移序列\footnote{译者注：原书在紧邻的尾能量式中
交替使用 $u_{i_j,\tau_j}$、$u_{i_j,\tau_{+,j}}$ 与 $u_{i_j,\tau_{i_j}}$；随后又把
这一轮新平移称为“序列 $v_k$”，却紧接着定义 $w_k$。选定
$w_k(\nu_k,t_k)$ 后，原书写 $\widetilde w_k=w_k\circ\nu_k$，这不能把该点移至
$0$；继而又把复合平移写成 $u_{i_k}\circ(-\mu_k+\nu_k)$，并遗漏原恢复平移
$\tau_{i_k}$。译文统一使用 $u_{i_j,\tau_{i_j}}$ 与 $w_k$，并依本节约定
$u\circ a=u(\mathord\cdot-a,\mathord\cdot)$ 改用 $w_k\circ(-\nu_k)$，再合并
三次平移。}
\[
  w_k:=u_{i_k,\tau_{i_k}}\circ\mu_k,
  \qquad
  w_k(\tau,t)=u_{i_k,\tau_{i_k}}(\tau-\mu_k,t).
\]
由定义及~\eqref{eq:half-tail-energy-cutoffs}，对所有 $k$ 都有
\[
  E_{J,(-\infty,0]}(w_k)
  >\frac16\left(
    \omega(B)-\omega(u_{(I,\tau_I)})
  \right).
\]
因此，可选择 $(\nu_k,t_k)\in(-\infty,0]\times[0,1]$，使得
\begin{equation}
  w_k(\nu_k,t_k)
  =u_{i_k,\tau_{i_k}}(\nu_k-\mu_k,t_k)
  \in M\setminus U,
\label{eq:new-recovery-point-outside-intersection}
\end{equation}
其中 $U$ 是 $L_0\cap L_1$ 的一个充分小邻域。考虑映射
\[
  \widetilde w_k
  :=w_k\circ(-\nu_k)
  =u_{i_k}\circ(\tau_{i_k}+\mu_k-\nu_k)
\]
在
\[
  (-\infty,\mu_k-R-\nu_k]\times[0,1]
\]
上的限制。依定义，$\widetilde w_k(0,t_k)\in M\setminus U$。注意，由于
$\nu_k\leq0$，仍有
\[
  \mu_k-R-\nu_k\longrightarrow\infty.
\]
由 $[0,1]$ 与 $M$ 的紧性，可在子列 $I'\subset\{i_k\}_{k\in\mathbb N}$ 上
同时假设
\[
  t_k\to t_\infty,
  \qquad
  \widetilde w_k(0,t_k)\to x.
\]
于是 $\widetilde w_k$ 有一个局部极限 $\widetilde w_\infty$，满足
\[
  \widetilde w_\infty(0,t_\infty)=x
  \in M\setminus(L_0\cap L_1).
\]
由于已知 $\widetilde w_\infty(\pm\infty)\in L_0\cap L_1$，它不可能是定常的。
相应的\termRecoveringPair{}记为\footnote{译者注：原书在这几行把
$\widetilde w_k$ 漏写下标为 $\widetilde w$，把局部极限先记为 $v'_\infty$、随后
又记为 $\widetilde w_\infty$，并把新恢复参数写成
$\tau'_\ell=\mu_\ell+\tau_\ell$，没有交代 $\tau_\ell$ 指原恢复平移还是 $\nu_k$。
译文统一映射字母，并依上面三次平移的实际复合写出新参数。}
$(I',\tau'_{I'})$，其中
\[
  \tau'_{i_k}=\tau_{i_k}+\mu_k-\nu_k,
  \qquad i_k\in I'.
\]
由构造，
\[
  [\widetilde w_\infty]
  <
  [(\widetilde w_\infty,u_{(I,\tau_I)})].
\]
如果
\[
  \omega([\widetilde w_\infty])
  +\omega([u_{(I,\tau_I)}])
  =\omega(B),
\]
那么 $(\widetilde w_\infty,u_{(I,\tau_I)})$ 又构成一条\termMaximalChain{}；否则可
重复上述过程，得到一条更大的链。由引理 14.2.12，这个过程必在有限步后停止，
命题得证。
\end{proof}

\begin{definition}
\label{def:gluable-floer-chain}
若\personFloer{}流线链 $(v_1,\ldots,v_N)$ 满足
\[
  v_i(\infty)=v_{i+1}(-\infty),
  \qquad i=1,\ldots,N-1,
\]
则称它是\termGluable{}的。
\footnote{译者注：原书把指标范围写成 $i=1,\ldots,N$，会在 $i=N$ 时引入
未定义的 $v_{N+1}$。译文补正为 $i=1,\ldots,N-1$。}
\end{definition}

下面陈述\termRecoveringSequence{}的\termMaximalChain{}及其相应\termLocalLimit{}的一个
基本性质。

\begin{proposition}
\label{prop:maximal-chain-consecutively-gluable}
设 $u_i\in\widetilde{\mathcal M}(p,q;B)$ 满足一致估计
\[
  \|du_i\|_\infty<C.
\]
对每条\termMaximalChain{}\footnote{译者注：原书在同一陈述中先把链长记为 $\ell$，
端点式却改用 $N$；端点处还把局部极限改写成未定义的 $u_{I_1},u_{I_N}$。译文
统一以 $N$ 表示链长，并写回首末两个局部极限。}
\[
  C_{(I,\vec\tau_I)}
  =
  \bigl(
    u_{(I,\tau_I^1)},\ldots,u_{(I,\tau_I^N)}
  \bigr)
\]
相应\termLocalLimit{}
\[
  \bigl(
    u_{(I,\tau_I^1)},\ldots,u_{(I,\tau_I^N)}
  \bigr)
\]
依次可粘合，并满足
\[
  u_{(I,\tau_I^1)}(-\infty)=x,
  \qquad
  u_{(I,\tau_I^N)}(\infty)=y.
\]
\end{proposition}

\begin{proof}
由命题 14.2.4，只需证明相邻的两个\termLocalLimit{}可粘合。把与它们相伴的\termRecoveringPair{}记为
\[
  \begin{aligned}
  I_i&=\{n_\ell^i\}_{\ell\in\mathbb N},
  &\tau_{I_i}&=\{\tau_\ell^i\}_{\ell\in\mathbb N},\\
  I_{i+1}&=\{n_\ell^{i+1}\}_{\ell\in\mathbb N},
  &\tau_{I_{i+1}}&=\{\tau_\ell^{i+1}\}_{\ell\in\mathbb N}.
  \end{aligned}
\]
反设存在 $i_0$，使相应的两个局部极限不可粘合。依定义，存在
$I''\subset I$，使得
\begin{equation}
  \tau_{\ell_k}^{i+1}-\tau_{\ell_k}^{i}
  \nearrow\infty
  \qquad(k\to\infty),
\label{eq:separating-recovery-parameters}
\end{equation}
其中 $I''=\{\ell_k\}_{k\in\mathbb N}$。现在考虑子列 $u_{\ell_k}$，记平均值
\begin{equation}
  \mu'_k
  =
  \frac{\tau_{\ell_k}^{i+1}+\tau_{\ell_k}^{i}}2,
\label{eq:midpoint-recovery-parameter}
\end{equation}
并考虑平移序列
\[
  u''_k=u_{\ell_k}\circ\mu'_k.
\]
可从中抽取一个有局部极限 $u''_\infty$ 的子列。把相应的\termRecoveringPair{}记为
$(I^{(3)},\tau_{I^{(3)}})$，其中 $I^{(3)}\subset I''$。由
\eqref{eq:separating-recovery-parameters}、\eqref{eq:midpoint-recovery-parameter}
中 $\mu'_k$ 的选取，有
\[
  (I_{i_0},\tau_{I_{i_0}})
  <(I^{(3)},\tau_{I^{(3)}})
  <(I_{i_0+1},\tau_{I_{i_0+1}}).
\]
这里的次序按相应局部极限所给出的链次序理解。
\footnote{译者注：原书在证明开头把只对\personFloer{}流线链定义的“可粘合”直接用于
恢复对；译文写回相邻局部极限。这里又把定义 14.2.11 只对链等价类定义的符号
$<$ 直接用于恢复对；译文保留所示比较式，并明示它经由相应局部极限的链次序
理解。}

这说明，除非 $(I^{(3)},\tau_{I^{(3)}})$ 对应的局部极限 $u^{(3)}$ 是定常的，否则
将与 $C=C_{(I,\vec\tau_I)}$ 为\termMaximalChain{}的假设矛盾。这便证明相邻局部
极限可粘合。类似的证明同时给出
$v_1(-\infty)=x$ 与 $v_N(\infty)=y$，从而完成证明。
\footnote{译者注：原书最后一句在两个都属于命题结论的端点等式之间使用 “or”。
译文依命题陈述补正为“与”。}
\end{proof}

下面证明，任意\termMaximalChain{}都恢复原同伦类。

\begin{theorem}
\label{thm:maximal-chain-recovers-homotopy-class}
设 $u_i\in\widetilde{\mathcal M}(p,q;B)$ 是任意满足
\[
  \|du_i\|_\infty<C
\]
的序列，并设
\[
  (v^1,\ldots,v^N)
\]
是 $\{u_i\}$ 的一条\termLocalLimitChain{}，对应于\termRecoveringSequence{}的一条\termMaximalChain{}。则
\begin{equation}
  B=[v^1]\#\cdots\#[v^N].
\label{eq:maximal-chain-homotopy-class}
\end{equation}
\end{theorem}

\begin{proof}
在证明引理 14.2.12 与命题 14.2.13 的过程中，我们已经得到
\[
  \omega(B)
  =
  \sum_{n=1}^N\int(v^n)^*\omega.
\]
这蕴含存在序列
\[
  \tau_{1,i}<\tau_{2,i}<\cdots<\tau_{N,i},
\]
使得
\begin{enumerate}
  \item $\tau_{\ell+1,i}-\tau_{\ell,i}\to\infty$（$i\to\infty$）；
  \item 对 $\ell=1,\ldots,N$，存在 $R_{\ell,i}>0$，满足
  $R_{\ell,i}\to\infty$，并且对 $\ell=1,\ldots,N-1$ 有
  \[
    R_{\ell,i}+R_{\ell+1,i}
    <\tau_{\ell+1,i}-\tau_{\ell,i},
  \]
  而且\footnote{译者注：原书先把 $R_{\ell,i}$ 的指标范围写成无上界的
  “$i=1,\ldots$”，译文依当前链长补正为 $\ell=1,\ldots,N$。公式 (14.2.35)
  又对每个局部极限中心附近的区间
  $[-R_{\ell,i},R_{\ell,i}]$ 积分并声称趋于零；这与 $u_i\circ\tau_{\ell,i}$
  收敛到非定常的 $v^\ell$ 直接矛盾。紧接着的句子实际把 $\epsilon$-正则性用于各
  局部极限之间的颈区。译文依该句把 (14.2.35) 的积分域补正为颈区之并。原书所列
  并集又把 $\ell$ 取到 $N$，引入未定义的 $\tau_{N+1,i}$，并在右端重复使用
  $R_{\ell,i}$；译文相应补正为 $\ell=1,\ldots,N-1$ 与 $R_{\ell+1,i}$。}
  \begin{equation}
    \int_{\displaystyle
      \bigcup_{\ell=1}^{N-1}
      [\tau_{\ell,i}+R_{\ell,i},
       \tau_{\ell+1,i}-R_{\ell+1,i}]\times[0,1]}
      u_i^*\omega
    \longrightarrow0
    \qquad(i\to\infty).
  \label{eq:vanishing-neck-area}
  \end{equation}
\end{enumerate}
一旦有~\eqref{eq:vanishing-neck-area}，\termEpsilonRegularityTheorem{}就蕴含 $u_i$
在区域
\[
  \bigcup_{\ell=1}^{N-1}
  [\tau_{\ell,i}+R_{\ell,i},
   \tau_{\ell+1,i}-R_{\ell+1,i}]\times[0,1]
\]
上的像包含于 $L_0\cap L_1$ 的一个小邻域中。利用 $L_0\cap L_1$ 的孤立性作一个
简单的拓扑论证，便得到
$B=[v^1]\#\cdots\#[v^N]$，证明完成。
\end{proof}

考虑从 $x$ 出发并在 $y$ 终止的可粘合链之等价类集合。我们将在
$\mathcal M(x,y;B)$ 的这个\termSplittingEnd{}中每一点处给出一个\termNeighborhoodBasis{}。
回忆
\[
  \mathcal M(x,y;B)
  =\widetilde{\mathcal M}(x,y;B)/\mathbb R
\]
是商空间。对 $\mathcal M(x,y;B)$ 中的每个元素，取其等价类中唯一满足\termEnergyNormalization{}
\begin{equation}
  E_J\bigl(u|_{(-\infty,0]\times[0,1]}\bigr)
  =
  E_J\bigl(u|_{[0,\infty)\times[0,1]}\bigr)
\label{eq:energy-normalized-slice}
\end{equation}
的映射 $u$ 作为代表。沿 $\tau$-方向平移任意给定的 $u$，总能唯一实现这个
归一化条件。以 $\widetilde{\mathcal M}_0$ 表示
$\widetilde{\mathcal M}$ 中由满足该归一化条件的 $u$ 组成的子集，并称它为一个\termSlice{}。商映射在 $\widetilde{\mathcal M}_0$ 上的限制于是给出到
$\mathcal M$ 上的微分同胚。为在\termSplittingEnd{}的每一点处给出上述\termNeighborhoodBasis{}，我们把这个\termSlice{}的构造推广到可粘合链。

\begin{definition}
\label{def:convergence-to-gluable-chain}
考虑一条映射的可粘合链
\[
  v^\ell\in
  \widetilde{\mathcal M}_0(x_{\ell-1},x_\ell;B_\ell),
\]
即
\[
  (v^1,\ldots,v^N),
  \qquad
  v^1(-\infty)=x,
  \quad
  v^N(\infty)=y,
\]
且 $B_1\#\cdots\#B_N=B$。若存在序列
$(\tau_{1,i},\ldots,\tau_{N,i})$，满足
\[
  \tau_{1,i}<\cdots<\tau_{N,i},
\]
并使
\[
  (v^1,\ldots,v^N)
  =C_{(\mathbb N,\vec\tau_{\mathbb N})}
\]
是序列 $u_i$ 的一条\termMaximalChain{}，则称序列
$\{[u_i]\}\subset\mathcal M(x,y;B)$ 收敛到
$[v^1]\#\cdots\#[v^N]$。等价地，向量
\[
  \vec\tau_{\mathbb N,i}
  =(\tau_{1,i},\ldots,\tau_{N,i})
\]
满足定理 14.2.17 的证明中所述性质 (1)、(2)。
\end{definition}

也可以在 $(v^1,\ldots,v^N)$ 处，为 $\mathcal M(x,y;B)$ 构造如下无穷远\termNeighborhoodBasis{}。取充分大的正数 $R_1,\ldots,R_N$，考虑
\[
  \widetilde{\mathcal V}^{\epsilon}_{\delta}
  \bigl((v^1,\ldots,v^N)\bigr)
\]
这一族 $\mathcal F(x,y;B)$ 的开子集\footnote{译者注：原书没有量化 $R_\ell$，并在
定义一个由映射 $u$ 组成的开集时突然改用 $u_i,\tau_{\ell,i}$。公式 (14.2.37) 又在
已经平移的 $u_i\circ\tau_{\ell,i}$ 上使用仍以 $\tau_\ell$ 为中心的区间；
公式 (14.2.38) 同样混合平移前后的坐标，并声称对 $\ell=1,\ldots,N$ 成立，从而
引入未定义的 $\tau_{N+1}$。译文按前述链收敛定义统一为单个映射 $u$：在平移后的
固定窗口上量 $W^{1,p}$ 距离，在未平移的相邻中心之间量颈能量，并把后者的范围
补正为 $\ell=1,\ldots,N-1$。原书本节没有定义 $\mathcal F(x,y;B)$，译文保留
该环境空间记号。}；其中每个子集由满足以下条件的映射 $u$
组成：存在 $\tau_1<\cdots<\tau_N$，使
\begin{align}
  d_{W^{1,p}([-R_\ell,R_\ell]\times[0,1])}
  (u\circ\tau_\ell,v^\ell)
  &<\epsilon,
  &&\ell=1,\ldots,N,
  \label{eq:gluable-chain-neighborhood-distance}\\
  E_J\bigl(
    u|_{[\tau_\ell+R_\ell,
           \tau_{\ell+1}-R_{\ell+1}]\times[0,1]}
  \bigr)
  &<\delta,
  &&\ell=1,\ldots,N-1.
  \label{eq:gluable-chain-neighborhood-neck-energy}
\end{align}
注意，$\widetilde{\mathcal V}^{\epsilon}_{\delta}
((v^1,\ldots,v^N))$ 在定义域的 $\tau$-平移下不变。

由上述\termNeighborhoodBasis{}的定义与\termExponentialDecay{}，立即得到下面的
引理。

\begin{lemma}
\label{lem:neighborhood-characterization-chain-convergence}
序列 $[u_i]$ 依定义 14.2.18 的意义收敛到
$[v^1]\#\cdots\#[v^N]$，当且仅当\footnote{译者注：原书把“对每个
$\epsilon,\delta>0$”写成“存在某些 $\epsilon,\delta>0$”。最终落入某一个固定
邻域并不刻画拓扑收敛；依\termNeighborhoodBasis{}的定义，必须最终落入每一个
这样的邻域。译文补正量词。}对每个 $\epsilon,\delta>0$，都存在
$N_0\in\mathbb N$，使得对所有 $i\geq N_0$，都有
\[
  u_i\in
  \widetilde{\mathcal V}^{\epsilon}_{\delta}
  \bigl((v^1,\ldots,v^N)\bigr).
\]
\end{lemma}

\begin{exercise}
\label{ex:neighborhood-characterization-chain-convergence}
证明这个引理。
\end{exercise}

因此，集合族
\begin{equation}
  \widetilde{\mathcal V}^{\epsilon}_{\delta}
  \bigl((v^1,\ldots,v^N)\bigr)
  \cap
  \widetilde{\mathcal M}(x,y;B)/\mathbb R
  =:
  \mathcal V^{\epsilon}_{\delta}
  \bigl((v^1,\ldots,v^N)\bigr)
\label{eq:splitting-end-neighborhood-basis}
\end{equation}
在 $[v^1]\#\cdots\#[v^N]$ 处给出 $\mathcal M(x,y;B)$ 的一个无穷远\termNeighborhoodBasis{}。

到目前为止，我们一直假设序列 $u_i\in\mathcal M(p,q;B)$ 具有 $C^1$ 界。一般说来，
当一对拉格朗日子流形 $(L_0,L_1)$ 周围存在非平凡全纯圆盘或球面时，这个 $C^1$
界并不成立。不过，仍可应用\termSacksUhlenbeckBubblingArgument{}，并证明在收敛过程
中移去气泡后，同一类收敛依然成立。
\footnote{译者注：原书把同一论证的作者姓氏印成 “Sachs--Uhlenbeck”；第一卷及
文献 (SU81) 均为 “Sacks--Uhlenbeck”。译文依跨卷既有裁决采用正确拼法。}

下一节将精确描述为紧化 $\mathcal M(p,q;B)$ 所要加入的对象。
